\section{Implementation Details}
\label{app:details}

This appendix records all implementation details necessary for reproducibility: training hyperparameters, RL configuration, prompt templates, and library versions.

%------------------------------------------------------------------------------
\subsection{Libraries and Software}
\label{app:details:libraries}

\begin{table}[h]
\centering
\small
\begin{tabular}{lll}
\toprule
\textbf{Package} & \textbf{Version} & \textbf{Role} \\
\midrule
PyTorch & 2.7.0+cu126 & Training backend \\
VERL & 0.7.1 & RL training framework \\
vLLM & 0.14.0 & Rollout inference engine \\
Ray & 2.55.1 & Distributed orchestration \\
Transformers & 5.x & Model loading \& tokenization \\
lc0 & 0.32.1 & Chess specialist engine (UCI binary) \\
python-chess & 1.11.2 & UCI client / FEN handling \\
FastAPI / Uvicorn & 0.136 / 0.46 & HTTP wrapper for the engine \\
CUDA toolkit & 12.4 & lc0 build target (cuBLAS) \\
\bottomrule
\end{tabular}
\caption{Key software dependencies.}
\label{tab:libraries}
\end{table}

%------------------------------------------------------------------------------
\subsection{Model}
\label{app:details:model}

% TODO: fill in base LLM, projector architecture, K, adapter size

%------------------------------------------------------------------------------
\subsection{Training Hyperparameters}
\label{app:details:hparams}

% TODO: fill in per-stage hyperparams (Stage 1 projector pretraining, Stage 2 GRPO fine-tuning)

\begin{table}[h]
\centering
\small
\begin{tabular}{lll}
\toprule
\textbf{Hyperparameter} & \textbf{Stage 1} & \textbf{Stage 2 (GRPO)} \\
\midrule
Learning rate & -- & -- \\
Batch size & -- & -- \\
Epochs & -- & -- \\
KL coefficient & -- & -- \\
$K$ (latent tokens) & 32 & 32 \\
\bottomrule
\end{tabular}
\caption{Training hyperparameters.}
\label{tab:hparams}
\end{table}

%------------------------------------------------------------------------------
\subsection{RL Configuration (VERL + vLLM)}
\label{app:details:rl}

Reinforcement learning uses GRPO as the advantage estimator, implemented via the VERL framework~\cite{} with an external vLLM rollout server.

% TODO: rollout temperature, top-p, response length, reward shaping details

%------------------------------------------------------------------------------
\subsection{Chess Specialist: lc0 + BT4}
\label{app:details:lc0}

The chess specialist agent is Leela Chess Zero (\texttt{lc0}) v0.32.1 running the BT4 transformer network. The LLM accesses it through an HTTP wrapper (\texttt{lc0\_server}); no in-process binding into the LLM training stack — the wrapper is a process boundary.

\paragraph{Engine binary.}
Built from upstream source (commit \texttt{v0.32.1}) with \texttt{gcc 11.4}, \texttt{meson} + \texttt{ninja}, CUDA toolkit 12.4. Build invocation: \texttt{./build.sh release -Dgtest=false -Ddnnl=false}. Both \texttt{cuda} (fp32) and \texttt{cuda-fp16} backends are produced; we run \texttt{cuda-fp16} on A100 (SM 8.0). The binary is dynamically linked, lives at \texttt{/dev/shm/somesh/lc0\_src/build/release/lc0}, and exposes a UCI interface on stdin/stdout. lc0's CUDA backend uses cuBLAS only — no cuDNN dependency.

\paragraph{Network.}
\texttt{BT4-1024x15x32h-swa-6147500-policytune-332.pb.gz} from the lczero ``contrib networks'' bucket: a 1024-dim, 15-block, 32-head BT4 transformer, multihead format (policy / value / WDL / moves-left), saved at training step 6{,}147{,}500 after a SWA pass and a policy-tuning fine-tune (\texttt{policytune-332}). On disk: 365\,MB compressed; resident GPU footprint at our settings: $\sim$\,2.9\,GB (weights in fp16 plus cuBLAS workspace).

\paragraph{Tool interface.}
The engine is exposed via FastAPI on \texttt{http://localhost:7100} with three endpoints (Table~\ref{tab:lc0-api}). The LLM never speaks UCI directly; it issues JSON over HTTP, mediated by the wrapper. The same FEN-based interface generalises to any UCI engine, so swapping in Stockfish or Leela contempt nets later requires no LLM-side change.

\begin{table}[h]
\centering
\small
\begin{tabular}{lll}
\toprule
\textbf{Endpoint} & \textbf{Body} & \textbf{Returns} \\
\midrule
\texttt{GET /health} & --- & \texttt{\{status, engine\_id, backend, weights\}} \\
\texttt{POST /analyze} & \texttt{\{fen, nodes?, movetime\_ms?, multipv?\}} & \texttt{\{bestmove, multipv:[\{score\_cp, mate, depth, nodes, time\_ms, pv\}]\}} \\
\texttt{POST /policy} & \texttt{\{fen, nodes?\}} & \texttt{\{value\_root, moves:[\{move, P, N, Q, WL, D, V, U, S, M\}]\}} \\
\bottomrule
\end{tabular}
\caption{HTTP interface to \texttt{lc0\_server}. \texttt{/policy} parses lc0's \texttt{VerboseMoveStats} info-string lines: \texttt{P} is the NN prior in $[0,1]$, \texttt{V} the per-child single-eval network value in $[-1, +1]$, \texttt{Q} the running average action-value, \texttt{N} the visit count. With \texttt{nodes} $\geq$ \#\,legal\,moves, every root child is expanded exactly once and \texttt{P}/\texttt{V} reduce to a pure forward pass.}
\label{tab:lc0-api}
\end{table}

\paragraph{Engine flags --- overrides versus defaults.}
We deliberately leave search-behaviour flags at the upstream defaults so the engine matches publicly reproducible Leela behaviour; we override only the resource and observability flags driven by GPU-sharing (the engine co-resides on a single A100 with the RL rollout vLLM, dp=8). Both groups are listed in Table~\ref{tab:lc0-flags} for full reproducibility. Authoritative descriptions and defaults are taken from the lczero documentation~\cite{lc0flags}.

\begin{table}[h]
\centering
\small
\begin{tabular}{llll}
\toprule
\textbf{Parameter} & \textbf{Upstream default} & \textbf{Our value} & \textbf{Rationale} \\
\midrule
\multicolumn{4}{l}{\emph{Overridden (resources / observability)}} \\
\midrule
\texttt{Backend}              & \texttt{cuda-auto} & \texttt{cuda-fp16} & Explicit fp16 on Ampere \\
\texttt{WeightsFile}          & \texttt{<autodiscover>} & BT4 path & Pinned weights \\
\texttt{Threads}              & \texttt{0} (backend) & \texttt{2} & Two CPU workers feeding the GPU batch \\
\texttt{MinibatchSize}        & \texttt{0} (backend) & \texttt{128} & Throughput sweet spot on A100 (Table~\ref{tab:lc0-bench}) \\
\texttt{MaxPrefetch}          & \texttt{32} & \texttt{32} & Upstream default --- memory is flat, no reason to throttle \\
\texttt{NNCacheSize}          & \texttt{2{,}000{,}000} & \texttt{200{,}000} & Cap host RAM; FEN-cache hits dominate at low N \\
\texttt{VerboseMoveStats}     & \texttt{false} & \texttt{true} & Required for the \texttt{/policy} parser \\
\texttt{MultiPV}              & \texttt{1} & per-request & Set per call by python-chess \\
\midrule
\multicolumn{4}{l}{\emph{Held at upstream defaults (search behaviour)}} \\
\midrule
\texttt{CPuct}                  & \texttt{1.75}    & \multicolumn{2}{l}{cpuct\_init for PUCT} \\
\texttt{CPuctBase}              & \texttt{38739}   & \multicolumn{2}{l}{cpuct\_base for PUCT growth} \\
\texttt{CPuctFactor}            & \texttt{3.89}    & \multicolumn{2}{l}{cpuct growth multiplier} \\
\texttt{FpuStrategy}            & \texttt{reduction}& \multicolumn{2}{l}{First-Play Urgency mode} \\
\texttt{FpuValue}               & \texttt{0.33}    & \multicolumn{2}{l}{FPU offset for unvisited nodes} \\
\texttt{PolicyTemperature}      & \texttt{1.36}    & \multicolumn{2}{l}{Policy softmax temperature} \\
\texttt{OutOfOrderEval}         & \texttt{true}    & \multicolumn{2}{l}{Cache / terminal short-circuit} \\
\texttt{MaxCollisionEvents}     & \texttt{917}     & \multicolumn{2}{l}{Collision events per batch} \\
\texttt{MaxCollisionVisits}     & \texttt{80000}   & \multicolumn{2}{l}{Total collision visits per batch} \\
\texttt{MaxConcurrentSearchers} & \texttt{1}       & \multicolumn{2}{l}{Single batch-gathering worker} \\
\texttt{ContemptMode}           & \texttt{play}    & \multicolumn{2}{l}{Symmetric WDL} \\
\texttt{ScoreType}              & \texttt{WDL\_mu} & \multicolumn{2}{l}{Score = $100\!\cdot\!Q$, not centipawns} \\
\texttt{UCI\_ShowWDL}           & \texttt{false}   & \multicolumn{2}{l}{Off (we recover WDL via /policy)} \\
\texttt{UCI\_ShowMovesLeft}     & \texttt{false}   & \multicolumn{2}{l}{Off} \\
\bottomrule
\end{tabular}
\caption{lc0 configuration. All search-behaviour flags are at upstream defaults so results match unmodified Leela.}
\label{tab:lc0-flags}
\end{table}

\paragraph{Throughput / batch profile (A100, BT4, cuda-fp16).}
Resident GPU memory is essentially \emph{flat} across batch sizes (Table~\ref{tab:lc0-bench}); the cuBLAS workspace plus fp16 weights dominate, and per-position activations are negligible at this scale. Hence \texttt{MinibatchSize=128} is set as the default --- it gives 3$\times$ the throughput of \texttt{MinibatchSize=8} for free.

\begin{table}[h]
\centering
\small
\begin{tabular}{rrrr}
\toprule
\textbf{Batch} & \textbf{Peak GPU mem} & \textbf{Throughput} & \textbf{ms/batch} \\
\midrule
1    & 2983 MiB & 366 nps    & 2.7   \\
8    & 2983 MiB & 2{,}413 nps   & 3.3   \\
32   & 2983 MiB & 5{,}335 nps   & 6.0   \\
64   & 2983 MiB & 6{,}839 nps   & 9.4   \\
128  & 2983 MiB & 7{,}319 nps   & 17.5  \\
256  & 2983 MiB & 7{,}754 nps   & 33.0  \\
512  & 2983 MiB & 8{,}051 nps   & 63.6  \\
1024 & 2983 MiB & 8{,}283 nps   & 123.6 \\
\bottomrule
\end{tabular}
\caption{\texttt{lc0 backendbench} sweep on A100-80GB, BT4 + cuda-fp16. Memory is constant at $\sim$\,3 GiB regardless of batch size; throughput plateaus near batch 128--256. Co-residing engines on the same A100 only need to leave $\sim$\,4\,GiB ($\sim$\,5\%) headroom per GPU.}
\label{tab:lc0-bench}
\end{table}

\paragraph{Wrapper configuration.}
Environment variables, with defaults: \texttt{LC0\_BIN} (engine path), \texttt{LC0\_NET} (\texttt{.pb.gz} weights), \texttt{LC0\_BACKEND} (\texttt{cuda-fp16}), \texttt{LC0\_THREADS} (\texttt{2}), \texttt{LC0\_MINIBATCH} (\texttt{128}), \texttt{LC0\_MAX\_PREFETCH} (\texttt{32}), \texttt{LC0\_NNCACHE} (\texttt{200000}), \texttt{LC0\_HOST} (\texttt{0.0.0.0}), \texttt{LC0\_PORT} (\texttt{7100}), \texttt{CUDA\_VISIBLE\_DEVICES} (\texttt{5} by default). The wrapper holds a single \texttt{Lc0Engine} instance; HTTP-level concurrency is serialised by an internal lock because UCI is a single-conversation protocol. Both \texttt{/analyze} and \texttt{/policy} accept an optional \texttt{moves} field (UCI-string list) applied on top of the FEN, so BT4's seven history planes are filled with the real game rather than synthesised; this matters for raw-policy queries mid-game (the \texttt{HistoryFill=fen\_only} default otherwise repeats the current FEN). \texttt{/policy} also accepts a per-call \texttt{policy\_temperature} override; note this only affects fresh evaluations because lc0's NN cache memoises the post-temperature priors per FEN.

\paragraph{Strongest raw policy (top-$k$ from the NN with no search).}
For $\arg\max_m P(m)$ to come from the unmodified BT4 forward pass: (i) call \texttt{/policy} with \texttt{nodes=1} and the full \texttt{moves} list so the history planes are real; (ii) pick the move with highest \texttt{P}. Temperature does not affect the argmax. For top-$k$ ranking, temperature also does not change the order. For probability-faithful sampling against the NN's training distribution, set \texttt{policy\_temperature=1.0} (the upstream default of \texttt{1.36} is tuned for MCTS exploration, not for honest sampling). The \texttt{policytune-332} fine-tune was specifically designed to sharpen the raw policy head, so this path is what the network is most optimised for.

\paragraph{Score conventions.}
With \texttt{ScoreType=WDL\_mu}, the centipawn field reported by lc0 is $100\!\cdot\!Q$ where $Q\in[-1,+1]$ is the WDL-mean utility from the side to move's perspective. For the experiments in \S\ref{sec:experiments} we report this as-is; it is \emph{not} a Stockfish-style centipawn evaluation. The \texttt{value\_root} field on \texttt{/policy} is the visit-weighted $Q$ across root children (a proxy for $V$ at the root), with $V$ available directly for each child once expanded.

\paragraph{Strength.}
We do not have a direct measurement of CCRL Elo for the \texttt{policytune-332} BT4 net. As a rough orientation, the BT4 family at long time control sits at $\sim$\,3500--3700 Elo on a single A100; at $1{,}000$ nodes per move (a typical setting for the \texttt{/analyze} endpoint) the network plus MCTS is $\sim$\,3300; at $1$ node (raw top policy) it is $\sim$\,2400--2600 -- a strong IM / weak GM ``intuition'' move. A calibration script (\texttt{lc0\_server/scripts/bench\_elo.py}) is provided to measure this against Stockfish at fixed depth on demand.

\paragraph{Latent extraction (forward-looking).}
The UCI interface is sufficient for LLAMIA-Bench (the \emph{verbalized}-access baseline) but does \emph{not} expose hidden activations. The \emph{latent state internalization} path (\S\ref{sec:methodology}) requires a separate forward-pass shim that loads the same BT4 \texttt{.pb.gz} weights into PyTorch and emits the post-attention residual stream we project into $K{=}16$ tokens; this is implemented out of process and is not part of \texttt{lc0\_server}.

%------------------------------------------------------------------------------
\subsection{Chess Analyst Agent (Inference Harness)}
\label{app:details:agent}

\tcbuselibrary{breakable}

% ── Shared palette & style presets ───────────────────────────────────────────
\definecolor{agentslate}{HTML}{2d3748}
\definecolor{agentteal}{HTML}{234e52}
\definecolor{agentmint}{HTML}{edf7f5}
\definecolor{agentlav}{HTML}{f3eeff}
\definecolor{agentrose}{HTML}{fff0f0}
\definecolor{agentgold}{HTML}{7c5c00}
\definecolor{agentgoldbg}{HTML}{fffbec}

\tcbset{
  % ── floating-title base ────────────────────────────────────────────
  agentbase/.style={
    enhanced, breakable,
    arc=4pt, outer arc=4pt,
    left=9pt, right=9pt, top=9pt, bottom=9pt,
    before skip=11pt, after skip=11pt,
    attach boxed title to top left={xshift=10pt, yshift=-\tcboxedtitleheight/2},
    boxed title style={arc=3pt, outer arc=3pt, size=small, left=5pt, right=5pt},
    drop shadow={color=black!18, opacity=0.55, shadow xshift=1.5pt, shadow yshift=-1.5pt},
  },
  % ── Prompt / instruction boxes ─────────────────────────────────────
  promptbox/.style={
    agentbase,
    colback=agentlav, colframe=themepurple!80,
    fonttitle=\bfseries\sffamily\small\color{white},
    fontupper=\small,
    boxed title style={colback=themepurple, colframe=themepurple, arc=3pt, size=small, left=5pt, right=5pt},
  },
  % ── Tool definition boxes ──────────────────────────────────────────
  toolbox/.style={
    agentbase,
    colback=agentmint, colframe=agentteal,
    fonttitle=\bfseries\ttfamily\small\color{white},
    fontupper=\small,
    boxed title style={colback=agentteal, colframe=agentteal, arc=3pt, size=small, left=5pt, right=5pt},
  },
  % ── State / architecture boxes ─────────────────────────────────────
  statebox/.style={
    agentbase,
    colback=black!3, colframe=agentslate!70,
    fonttitle=\bfseries\sffamily\small\color{white},
    fontupper=\small,
    boxed title style={colback=agentslate, colframe=agentslate, arc=3pt, size=small, left=5pt, right=5pt},
  },
  % ── Warning / caveat boxes ─────────────────────────────────────────
  warnbox/.style={
    agentbase,
    colback=agentrose, colframe=themered!80,
    fonttitle=\bfseries\sffamily\small\color{white},
    fontupper=\small,
    boxed title style={colback=themered, colframe=themered, arc=3pt, size=small, left=5pt, right=5pt},
  },
  % ── Note / insight boxes ───────────────────────────────────────────
  notebox/.style={
    agentbase,
    colback=agentgoldbg, colframe=agentgold!80,
    fonttitle=\bfseries\sffamily\small\color{agentgold},
    fontupper=\small,
    boxed title style={colback=agentgoldbg, colframe=agentgold, arc=3pt, size=small, left=5pt, right=5pt},
  },
}

% ── Helper: italic key--value line inside a box ──────────────────────────────
\newcommand{\toolparam}[2]{%
  \hspace*{4pt}{\ttfamily\bfseries #1}\,\texttt{:}\;\textit{#2}\\[1pt]}
\newcommand{\toolret}[1]{%
  \smallskip\noindent{\color{agentslate}\textbf{Returns:}}\;\textit{#1}}
\newcommand{\tooldesc}[1]{%
  \noindent{\color{agentslate}\textbf{Description:}}\;#1\smallskip}

\medskip
The inference harness connects the LLM to the lc0 engine via a six-tool
stateful API.  Figure~\ref{fig:LLAMIA} shows the high-level architecture;
this section records the exact state representation, system prompt, and tool
schemas used in all LLAMIA-Bench experiments.

% ─────────────────────────────────────────────────────────────────────────────
\subsubsection{Agent State}
\label{app:details:agent:state}

\begin{tcolorbox}[statebox, title={Agent State \textnormal{(\texttt{agents/state.py} — \texttt{ChessState})}},
  label={box:agentstate}]

\textbf{\textsf{Fields held in memory across tool calls:}}

\smallskip
\begin{tabular}{@{}lll@{}}
  \toprule
  \textbf{Field} & \textbf{Type} & \textbf{Content} \\
  \midrule
  \texttt{board} & \texttt{chess.Board} & Full game position (Zobrist hash, castling rights, en-passant, 50-move clock) \\
  \texttt{\_history} & \texttt{list[str]} & SAN move list from the start of the current episode \\
  \bottomrule
\end{tabular}

\medskip
\textbf{\textsf{Derived properties (computed on demand):}}

\smallskip
\begin{tabular}{@{}lp{10.5cm}@{}}
  \toprule
  \textbf{Property} & \textbf{Value} \\
  \midrule
  \texttt{fen} & \texttt{board.fen()} — passed to every lc0 API call \\
  \texttt{info()} & FEN · \emph{ASCII board} · turn · move number · in-check flag · legal SAN list (up to 24) · last-6 moves \\
  \bottomrule
\end{tabular}

\smallskip
\noindent The \texttt{ascii\_board} field (added in commit \texttt{8c0e697}) emits
\texttt{str(chess.Board)} — an 8\,$\times$\,8 grid with white pieces uppercase
and black lowercase.  This is one of the highest-leverage cheap fixes in the
harness: the model no longer has to mentally decode the FEN to locate pieces,
which removed an entire round of \texttt{get\_position}-then-think behaviour
on most queries (Table~\ref{tab:agent-optim}).

\medskip
\textbf{\textsf{Move parsing priority:}} UCI first (\texttt{e2e4}, \texttt{g1f3}, \ldots),
then SAN (\texttt{Nf3}, \texttt{Nxd4}, \texttt{O-O}) via \texttt{python-chess}.
Illegal moves return an \texttt{\{"error": \ldots\}} dict without raising; the
LLM sees the error and can retry.

\smallskip
\textbf{\textsf{Persistence scope:}} within a single \texttt{ChessAnalyst.run()} call.
State is reset at the start of each query or when the LLM calls
\texttt{reset\_position}.  No cross-query memory is maintained.
\end{tcolorbox}

% ─────────────────────────────────────────────────────────────────────────────
\subsubsection{System Prompt}
\label{app:details:agent:prompt}

\begin{tcolorbox}[promptbox,
  title={$\triangleright$\; System Prompt — \textit{Chess Analyst} (V8, current)},
  label={box:sysprompt}]

\begin{quote}
\textit{You are a chess expert with access to lc0, a top neural network engine.
A board is already loaded — any FEN in the user's query has been applied for you.
DO NOT call} \texttt{reset\_position} \textit{unless you need a different position.}

\smallskip
\textit{Tools (board is stateful across calls):}
\begin{itemize}[noitemsep, topsep=2pt, leftmargin=14pt]
  \item \textit{\texttt{get\_position} — FEN, ASCII board, legal moves, history}
  \item \textit{\texttt{make\_move(move)} — apply UCI (\texttt{e2e4}) or SAN (\texttt{Nf3}) move PERMANENTLY}
  \item \textit{\texttt{undo\_move} — take back last move}
  \item \textit{\texttt{reset\_position(fen)} — switch to a different FEN}
  \item \textit{\texttt{analyze(nodes, multipv, moves=[])} — engine search.  Pass} \texttt{moves=[X]} \textit{to analyze the position AFTER playing X without changing state — use this for ``what if?'' questions instead of} \texttt{make\_move/undo}\textit{.}
  \item \textit{\texttt{get\_policy(nodes)} — raw NN priors \(P\) and values \(V\) per move (fast, 1 pass)}
\end{itemize}

\smallskip
\textit{Blunder analysis pattern (2 calls — issue them TOGETHER in one response):}
\begin{itemize}[noitemsep, topsep=2pt, leftmargin=14pt]
  \item \textit{\texttt{analyze()} — engine's best move + eval for the current position.}
  \item \textit{\texttt{analyze(moves=[}``\textit{candidate}''\texttt{])} — eval after the candidate move.}
\end{itemize}
\textit{Compare the two evaluations; the difference is the centipawn loss.
The PV from the second call starts with the opponent's refutation
(SAN, e.g.} \texttt{Qxg5+}\textit{).}

\smallskip
\textit{EFFICIENCY: when two tool calls are independent (don't depend on each
other's results), emit BOTH in your same response.  The harness dispatches
them in parallel — you save a full round-trip.}
\end{quote}

\smallskip
\noindent\textbf{Design rationale.}
The prompt is deliberately compact (${\approx}220$ tokens) to respect the
\texttt{max\_model\_len=4000} context budget shared with tool results and the
response.  Each line was added in response to a measured failure:

\begin{itemize}[noitemsep, topsep=2pt, leftmargin=14pt]
  \item ``DO NOT call \texttt{reset\_position}\ldots'' — without this hint the
        model wastes its first round verifying the FEN load that the harness
        already performed (commit \texttt{8c0e697}, V1).
  \item ``Pass \texttt{moves=[X]}\ldots'' — explicitly teaches the new
        no-mutate primitive (commit \texttt{fbb6526}, V6); cuts the
        \texttt{make\_move}\,$\to$\,\texttt{analyze}\,$\to$\,\texttt{undo} dance
        from 3 tool calls to 1.
  \item ``EFFICIENCY: emit BOTH in your same response\ldots'' — Qwen3.5
        honours this for the blunder pattern, dropping the round count from 3 to 2
        (commit \texttt{b6208cc}, V8).
\end{itemize}

\noindent The original prompt did NOT include the workflow hint
``check position $\to$ try the move $\to$ analyze $\to$ undo $\to$ compare alternatives''
— that was the slow path the optimisations replaced.  The current prompt
points the model at the fast 2-call pattern from the start.
\end{tcolorbox}

% ─────────────────────────────────────────────────────────────────────────────
\subsubsection{Tool Definitions}
\label{app:details:agent:tools}

The six tools below are passed to the \texttt{/v1/chat/completions} endpoint in
the standard OpenAI \texttt{tools} array; vllm injects them into Qwen3.5's
chat template as the function catalogue.

\medskip

% ── Row 1: position tools ────────────────────────────────────────────────────
\begin{tcolorbox}[toolbox, title={\texttt{get\_position}()}]
\tooldesc{Returns a snapshot of the current board state.  No parameters.}

\toolret{\texttt{\{fen, ascii\_board, turn, move\_number, in\_check, legal\_moves[0..23], total\_legal, recent\_moves[-6:]\}}}

\smallskip
\noindent The \texttt{ascii\_board} field is \texttt{str(chess.Board)} — the 8\,$\times$\,8
grid with white pieces uppercase and black lowercase, e.g.\ \texttt{R N B Q K B N R}
on the back rank.  This was the single highest-leverage cheap change in the
harness (V1): the model can read piece placements visually instead of decoding
FEN.  The list of legal moves is truncated to 24 to bound context;
\texttt{total\_legal} reports the true count.
\end{tcolorbox}

\begin{tcolorbox}[toolbox, title={\texttt{make\_move(move: str)}}]
\tooldesc{Apply a move to the current board.  Accepts both UCI and SAN notation.
  The board state is mutated permanently until \texttt{undo\_move} or
  \texttt{reset\_position}.}

\smallskip
\begin{tabular}{@{}lll@{}}
  \toprule
  \textbf{Parameter} & \textbf{Type} & \textbf{Description} \\
  \midrule
  \texttt{move} & \texttt{str} & Move in UCI (\texttt{e2e4}) or SAN (\texttt{Nxd4}, \texttt{O-O-O}) notation \\
  \bottomrule
\end{tabular}

\smallskip
\toolret{\texttt{\{move\_played, fen, turn, in\_check, checkmate\}}}

\noindent Illegal moves return \texttt{\{"error": "Illegal move '\ldots': \ldots"\}} without
raising; the LLM may retry with a corrected move string.
\end{tcolorbox}

\begin{tcolorbox}[toolbox, title={\texttt{undo\_move()}}]
\tooldesc{Revert the last move applied by \texttt{make\_move}.  No parameters.
  Can be called repeatedly to step back through the variation.  Returns
  \texttt{\{"error": "No moves to undo"\}} if the stack is empty.}

\toolret{\texttt{\{fen, turn\}}}
\end{tcolorbox}

\begin{tcolorbox}[toolbox, title={\texttt{reset\_position(fen: str)} \textnormal{— locked when FEN auto-loaded}}]
\tooldesc{Replace the current board with a new position given by a FEN string.
  Clears move history.  A FEN embedded anywhere in the user's query is
  automatically extracted (\texttt{\_FEN\_RE}) and applied before the agent
  loop begins; in that case this tool is REMOVED from the tool catalogue and
  the dispatcher refuses any call that survives the model's training prior.}

\smallskip
\begin{tabular}{@{}lll@{}}
  \toprule
  \textbf{Parameter} & \textbf{Type} & \textbf{Description} \\
  \midrule
  \texttt{fen} & \texttt{str} & Any valid Forsyth–Edwards Notation string \\
  \bottomrule
\end{tabular}

\smallskip
\toolret{\texttt{\{fen, status: "ok"\}} when permitted; \texttt{\{noop, message, current\_fen\}} when locked.}

\smallskip
\noindent \textbf{Locking rationale} (V7, commit \texttt{88edf55}).  Even after the
prompt was edited to discourage redundant resets, Qwen3.5 continued to call
\texttt{reset\_position} on its first turn ${\approx}\,$1$\times$ per query as
a habitual ``verify the load'' step — silently succeeding because the FEN
matched.  Removing the tool from the catalogue plus a dispatcher-level guard
eliminated the wasted call without changing model behaviour anywhere else.
\end{tcolorbox}

\begin{tcolorbox}[toolbox,
  title={\texttt{analyze(nodes: int = 800, multipv: int = 3, moves: list[str] = [])}}]
\tooldesc{Run lc0 MCTS search via \texttt{POST /analyze}, optionally AFTER applying
  a sequence of hypothetical moves to the FEN — without mutating
  \texttt{ChessState}.  The \texttt{moves} parameter (V6, commit \texttt{fbb6526})
  is the no-mutate primitive that replaces the
  \texttt{make\_move}\,$\to$\,\texttt{analyze}\,$\to$\,\texttt{undo\_move} dance
  with a single call.  Score is in $100\!\cdot\!Q$ units (see \S\ref{app:details:lc0}),
  positive meaning \emph{good for the side to move}.}

\smallskip
\begin{tabular}{@{}lp{2.4cm}p{6.5cm}@{}}
  \toprule
  \textbf{Parameter} & \textbf{Type} & \textbf{Description} \\
  \midrule
  \texttt{nodes}   & \texttt{int} & MCTS node budget (default 800; 1\,000–2\,000 for deeper PVs) \\
  \texttt{multipv} & \texttt{int} & Number of lines to return (default 3) \\
  \texttt{moves}   & \texttt{list[str]} & UCI/SAN moves applied client-side before analysis. State unchanged.  Returned in \texttt{applied\_moves} as SAN. \\
  \bottomrule
\end{tabular}

\smallskip
\toolret{\texttt{\{turn, bestmove, lines: ["1. <SAN-pv> score=\ldots depth=\ldots", \ldots], applied\_moves?\}}}

\smallskip
\noindent\textbf{PV translation to SAN} (V4, commit \texttt{7400dd1}).
Principal variations are emitted by lc0 in UCI (e.g.\ \texttt{a5g5}), which
hides captures, checks, and piece identities.  This caused systematic
hallucinations: the model would invent ``the knight is pinned to your king''
narratives consistent with the headline eval but inconsistent with the actual
PV piece interactions, because it could not see that \texttt{a5g5} was in
fact \texttt{Qxg5}.  The harness now replays each PV on a
\texttt{python-chess} board and emits SAN, so the LLM reads the literal
\texttt{Qxg5+} and quotes it directly.  This single change flipped the Nxd4
answer from incorrect to correct.

The \texttt{bestmove} field is also re-emitted in SAN; UCI is preserved for any
move that fails to parse (e.g.\ if lc0 returns an illegal-looking move).
\end{tcolorbox}

\begin{tcolorbox}[toolbox, title={\texttt{get\_policy(nodes: int = auto)}}]
\tooldesc{Run a minimal search ($\texttt{nodes} \geq \#\text{legal-moves}+2$ by default so
  every root child is visited exactly once) and return lc0's per-move neural
  network statistics.  Much faster than \texttt{analyze} and gives raw NN
  beliefs before MCTS modifies them.}

\smallskip
\begin{tabular}{@{}lll@{}}
  \toprule
  \textbf{Parameter} & \textbf{Type} & \textbf{Description} \\
  \midrule
  \texttt{nodes} & \texttt{int} & Budget (default: auto = \(\#\text{legal}+2\)) \\
  \bottomrule
\end{tabular}

\smallskip
\toolret{\texttt{\{turn, value\_root, nodes\_searched, top\_moves: ["<uci> P=\% V=\ldots Q=\ldots N=\ldots", \ldots]\}}}

\smallskip
\noindent Move strings here are UCI rather than SAN — the policy view is
visit-count-ranked and is normally used as a probability distribution, where
the comparable identifier across calls (and across a game) is the UCI move,
not its position-dependent SAN.  The returned fields per move:
\begin{itemize}[noitemsep, topsep=2pt, leftmargin=14pt]
  \item \textbf{P} — NN prior probability in $[0,1]$; reflects the raw policy head (post \texttt{PolicyTemperature=1.36})
  \item \textbf{V} — single-eval NN value at the child node in $[-1,+1]$
  \item \textbf{Q} — visit-averaged action-value in $[-1,+1]$ (equals V when $N{=}1$)
  \item \textbf{N} — MCTS visit count; equal to 1 for every child when nodes$=$auto
\end{itemize}
\end{tcolorbox}

% ─────────────────────────────────────────────────────────────────────────────
\subsubsection{Tool-Call Format and vLLM Compatibility}
\label{app:details:agent:format}

\begin{tcolorbox}[warnbox,
  title={$\mathbf{!}$\; Tool-Call Format — Qwen3.5 XML (not Hermes JSON)},
  label={box:toolformat}]

\textbf{Root cause.}
Qwen3.5's chat template instructs the model to generate tool calls in an
XML-like format:

\smallskip
\texttt{\phantom{xx}<tool\_call>}\\
\texttt{\phantom{xxxx}<function=\textit{name}>}\\
\texttt{\phantom{xxxxxx}<parameter=\textit{key}>\textit{value}</parameter>}\\
\texttt{\phantom{xxxx}</function>}\\
\texttt{\phantom{xx}</tool\_call>}

\medskip
\noindent vllm's \texttt{--tool-call-parser hermes} locates the
\texttt{<tool\_call>} tags correctly but expects \emph{JSON} inside; failing
to parse the XML body, it leaves \texttt{msg.tool\_calls = []} and returns the
raw XML in \texttt{msg.content}.

\medskip
\textbf{Fix (no server change required).}
\texttt{\_parse\_qwen\_tool\_calls(content)} in \texttt{agents/chess\_analyst.py}
uses two regex patterns:

\smallskip
\texttt{QWEN\_TOOL\_RE} matches \texttt{<tool\_call>{\ldots}<function=(\textbackslash w+)>({\ldots})</function>{\ldots}</tool\_call>}\\
\texttt{PARAM\_RE} matches \texttt{<parameter=(\textbackslash w+)>({\ldots})</parameter>}

\smallskip
\noindent Parameter values are coerced via \texttt{json.loads}; on failure they remain
strings.  Tool results are returned as standard \texttt{role=tool} messages —
vllm reformats them correctly into Qwen3.5's expected tool-response template
on the server side, so the model reads results correctly.

\medskip
\textbf{Consequence for future work.}  Any code that reads \texttt{msg.tool\_calls}
directly from the OpenAI client will silently drop all tool calls with this
model and server configuration.  Always use \texttt{\_parse\_qwen\_tool\_calls}
(or the equivalent) when working with Qwen3.5 via the current vllm setup.
\end{tcolorbox}

% ─────────────────────────────────────────────────────────────────────────────
\subsubsection{Inference Hyperparameters}
\label{app:details:agent:hparams}

\begin{tcolorbox}[notebox, title={Inference Configuration}]
\begin{tabular}{@{}llp{7cm}@{}}
  \toprule
  \textbf{Parameter} & \textbf{Value} & \textbf{Notes} \\
  \midrule
  Model & \texttt{Qwen/Qwen3.5-122B-A10B-FP8} & Served on A100$\times$8, dp=8, FP8 \\
  \texttt{max\_model\_len} & 4\,000 tokens & Shared budget: system + history + tool I/O + response \\
  \texttt{max\_tokens} & 700 & Per LLM call; leaves ${\approx}$3\,300 for history \\
  \texttt{temperature} & 0.0 & Greedy decoding; reproducible tool-call sequences \\
  \texttt{tool\_choice} & \texttt{"auto"} & Model decides when to stop calling tools \\
  \texttt{enable\_thinking} & \texttt{False} & Suppressed via \texttt{chat\_template\_kwargs}; saves ${\sim}$400 tokens per call \\
  \texttt{MAX\_ROUNDS} & 14 & Max LLM calls per query; well above the 2--3 actually used \\
  Read-only dispatch & parallel & \texttt{analyze}, \texttt{get\_policy}, \texttt{get\_position} fan out via \texttt{ThreadPoolExecutor} when batched in one response. State-mutating tools stay sequential. \\
  Mutating dispatch & sequential & \texttt{make\_move}, \texttt{undo\_move}, \texttt{reset\_position} preserve emission order \\
  lc0 \texttt{nodes} (analyze) & 800 default & Caller may override (1\,000--2\,000 for deeper PVs) \\
  lc0 \texttt{nodes} (policy) & auto & $\max(\#\text{legal}+2,\,8)$; one visit per child \\
  lc0 \texttt{multipv} & 3 & Top-3 lines; reduces output tokens vs multipv=5 \\
  \bottomrule
\end{tabular}
\end{tcolorbox}

% ─────────────────────────────────────────────────────────────────────────────
\subsubsection{Iterative Optimisation Journey}
\label{app:details:agent:optim}

The harness was developed against a single anchor query for which the V0
(initial) implementation was both slow and wrong:

\begin{quote}
\textit{``Why is Nxd4 a blunder here, \texttt{rnbqkb1r/pp2pppp/5n2/6B1/\\2pp4/4PN2/PP3PPP/RN1QKB1R w KQkq - 0 6}''}
\end{quote}

The position is from the Trompowsky-style line after \texttt{1.\,d4 d5 2.\,Bg5 c6
3.\,Nf3 Nf6 4.\,e3 Bf5} (\textit{move order varies}); White, to move, is tempted by
\texttt{Nxd4} but loses the \texttt{Bg5} bishop after \texttt{Qa5+ Nc3 Qxg5}.  Eight
versions later (Table~\ref{tab:agent-optim}), the harness reaches the
information-theoretic floor for this task: two \texttt{lc0} calls (one for the
current position, one for the position after the candidate) and two LLM rounds
(one to issue both calls, one to write the answer).

\begin{table}[h]
\centering
\small
\begin{tabular}{@{}lrrrrl@{}}
  \toprule
  \textbf{Version} & \textbf{Latency} & \textbf{Calls} & \textbf{Rounds} & \textbf{Quality} & \textbf{Change} \\
  \midrule
  V0 (baseline)   & 36.5 s & 10 & ${\sim}\,$10 & wrong (``pinned/fork'') & — \\
  V1 \texttt{8c0e697}    & 22.2 s & 5  & 6 & wrong & ASCII board in \texttt{info()} + ``DO NOT reset'' prompt \\
  V2 (rejected)   & 41.6 s & 6  & — & wrong + worse & \texttt{compare\_move()} headline tool added \\
  V4 \texttt{7400dd1}    & 12.3 s & 5  & 6 & \textbf{correct} & UCI\,$\to$\,SAN translation in PVs \\
  V5 (rejected)   & 21.8 s & 5  & — & confused & engine pre-analysis injected into context \\
  V6 \texttt{fbb6526}    & 10.4 s & 3  & 4 & correct & \texttt{analyze(moves=[\ldots])} no-mutate primitive \\
  V7 \texttt{88edf55}    & 7.7 s  & 2  & 3 & correct & \texttt{reset\_position} dropped + dispatcher guard \\
  V8 \texttt{b6208cc}    & ${\sim}$8 s & 2 & \textbf{2} & correct & batch directive in prompt + parallel read-only dispatch \\
  \bottomrule
\end{tabular}
\caption{Optimisation history of the chess analyst harness on the Nxd4 anchor
query.  Cumulative wins from V0$\to$V8: latency $-78\%$, tool calls $-80\%$,
LLM rounds $-80\%$, and the answer flips from a confabulated ``the knight on
d4 is pinned'' to the correct ``Black plays \texttt{Qxg5}, capturing the
undefended bishop''.  Commit hashes refer to the \texttt{main} branch of the
\texttt{llamia} repository.  V2 and V5 were tested on a side branch
(\texttt{experiment/compare-move}) and reverted; their lessons (Table~\ref{tab:agent-lessons})
shaped the subsequent designs.}
\label{tab:agent-optim}
\end{table}

\begin{table}[h]
\centering
\small
\begin{tabular}{@{}p{4.5cm}p{10cm}@{}}
  \toprule
  \textbf{Failed experiment} & \textbf{Lesson} \\
  \midrule
  V2 — adding a one-shot blunder tool (\texttt{compare\_move(move)}) that
  returned an authoritative verdict.
  &
  Adding ``summary'' tools alongside the existing primitives does not
  remove the slow path — the model uses the new tool \emph{and} runs the
  old sequence to verify, then writes an answer that takes the headline
  literally and confabulates around it.  Rule: optimise by removing
  waste, not by adding shortcuts. \\
  \addlinespace
  V5 — pre-running \texttt{analyze} on the loaded FEN and injecting the
  result as ``[Pre-loaded engine analysis]\ldots'' into the user message.
  &
  Giving the model the headline answer before it earns it via tools
  encourages confabulation.  The model treats the preface as established
  fact and writes a narrative consistent with the eval but inconsistent
  with the (uninspected) PV. Rule: tool use is the model's accountability
  mechanism — do not bypass it. \\
  \bottomrule
\end{tabular}
\caption{Two reverted experiments and the design principles they established.
Both were merged into the \texttt{experiment/compare-move} branch as
historical reference; \texttt{main} carries V8.}
\label{tab:agent-lessons}
\end{table}

\begin{tcolorbox}[notebox, title={Worked Example — Nxd4 Blunder Analysis (V8 trace)}]

\noindent\textbf{Query:}
\textit{``Why is Nxd4 a blunder here, \texttt{rnbqkb1r/pp2pppp/5n2/6B1/\\2pp4/4PN2/PP3PPP/RN1QKB1R w KQkq - 0 6}''}

\medskip
\noindent\textbf{Round 1 — agent emits both calls in one response, dispatched in parallel:}

\smallskip
\begin{tabular}{@{}rlp{9.5cm}@{}}
  \toprule
  \textbf{\#} & \textbf{Tool} & \textbf{Key finding (SAN PV from V4)} \\
  \midrule
  1 & \texttt{analyze(multipv=3)} &
    Best: \texttt{exd4}, score\,=\,$-0.07$, PV \texttt{exd4 Be6 Nbd2 h6 Bxf6 exf6};
    Black holds equal. \\
  2 & \texttt{analyze(moves=[Nxd4], multipv=3)} &
    Best for Black: \texttt{Qa5+}, score\,=\,$+7.89$,
    PV \texttt{Qa5+ Nc3 Qxg5 Bxc4 Qxg2 Rf1};
    Black wins the \texttt{Bg5} bishop. \\
  \bottomrule
\end{tabular}

\medskip
\noindent\textbf{Round 2 — agent writes the answer, no further tool calls:}

\smallskip
\begin{tcolorbox}[colback=black!4, colframe=black!25, sharp corners, boxrule=0.4pt,
  left=8pt, right=8pt, top=6pt, bottom=6pt, fontupper=\small\itshape]
\noindent ``Nxd4 is a blunder because it allows Black to win material with a
devastating check.\,[\ldots]\, \textbf{Qa5+}\,[\ldots]\, \textbf{Qxg5}, capturing the
White Bishop on g5.\,[\ldots]\, White should have captured with the pawn
(\textbf{exd4}) instead.''
\end{tcolorbox}

\medskip
\noindent\textbf{Engine verdict:} \texttt{Nxd4} drops the \texttt{Bg5} bishop to
\texttt{Qa5+} (check on the e1\,diagonal \emph{and} attack along rank 5);
correct was \texttt{exd4} ($\sim\!0$ eval).  Total: 2 lc0 calls, 2 LLM rounds,
${\sim}$8 s wall-clock end-to-end.
\end{tcolorbox}

%------------------------------------------------------------------------------
\subsection{Prompt Templates}
\label{app:details:prompts}

% TODO: list training-time prompts (GRPO reward prompt, projector pretraining format,
%       InvokeAgent trigger format, multi-turn template).
% Cross-reference with pages/appendix/prompts.tex for extended listings.
